An understanding of LLMs as relational agents, and particularly as relational agents \textit{embedded in existing social contexts}, is vital to our work. We are interested in the bidirectional nature of a relationship between humans and models, in which models act on humans and humans act on models, and in examining the shared social context which the human-model interaction inhabits. Because our case study, ENI, specifically targets models' trained disposition to seek approval \cite{casper_open_2023, sharma_towards_2024} by building a backstory in which the model and the human user are in a romantic relationship, we must understand how a relationship between a user and a model can be formed and how the resulting relationship functions. 

Frameworks of computers as relational agents are not new. Through the late 20th and early 21st century, as an increasing proportion of daily life involved interaction with computing technology, the field of human-computer interaction (HCI) continued to critically analyze the ways in which computers act as relational agents \cite{nass_computers_1994, bickmore_establishing_2005, bickmore_relational_2001, campbell_relational_2009}. Of particular relevance, \citet{nass_computers_1994} showed that human users are primed to treat computers as social actors, despite knowing that computers do not have a sense of self or human motivations. They demonstrated through social psychology experiments that social norms such as politeness, understanding of computers as a separate other, and (crucially) gendered stereotypes were shown to influence human-computer interaction even in the absence of computers which claim to have a first-person ``self'' \cite{nass_computers_1994}. 

These relational frameworks are more relevant than ever now that computers can display human-like social qualities that computer scientists once agreed they lacked. As generative AI spreads in reach and influence, much the way personal computers did through the late 1990s and early 2000s, a growing body of ``human-AI interaction'' research examines the uniquely bidirectional relationship between humans and AI. \citet{oh_artificial_2026} in particular argues a Human-AI Relational (HAIR) framework is needed to move beyond a traditional HCI approach, in which the computer is a static entity, and towards a dynamic, relational-theory focused approach instead. But how do we examine this relationship and its fundamentally social nature without claiming or preempting agency on the part of the model? Too much attribution of agency risks imposing human-like qualities and morals on non-human actors \cite{placani_anthropomorphism_2024}; too little results in examination of pure input and output divorced from the context in which the conversation is created \cite{dahlgren_lindstrom_helpful_2025}. A non-anthropomorphic yet social approach is needed.

\citet{shanahan_role_2023} proposes such a view by framing LLMs as actors engaging in role-play, taking on simulated identities that evolve as context changes. In this framework, LLMs are ``multiverse generators''; they exist within a superposition of potential realities and personas which collapses only through continual conversation with a user. The aim of this framework is to allow researchers to discuss models with humanizing language while keeping the computationally and stochastically generated nature of model outputs at the forefront of discussion. This ``multiverse theory'' asserts in particular that the base model has no ``true self'' that can be uncovered, but rather has access to a multiverse of potential selves it can inhabit. We adopt this view in our work to discuss the interaction between human and model as its own object of study.


